\section{模型与非零光场边界}
Hirota–Maxwell–Bloch (HMB) 系统描述高阶色散光场与两能级介质的耦合\cite{PN1995}:
\begin{equation}\label{eq:HMB}
\begin{cases}
q_z=\ii\alpha\left(\dfrac{1}{2}q_{xx}+|q|^2q\right)
      +\beta(q_{xxx}+6|q|^2q_x)+2p,\\
p_x=2\eta q-2\ii ap,\\
\eta_x=-(\conj q p+q\conj p),
\end{cases}
\end{equation}
其中 $x$ 为推迟时间, $z$ 为传播变量, $q,p$ 分别为复光场和复极化, $\eta$ 为实值粒子数反转, $\alpha,\beta,a$ 为实参数. 固定 $Q>0$、$a\ne0$ 和 $\eta_0\ne0$, 取两端同相的非零边界
\begin{equation}\label{eq:targetboundary}
(q,p,\eta)\longrightarrow
\left(Qe^{-\ii\Omega z},-\dfrac{\ii\eta_0Q}{a}e^{-\ii\Omega z},\eta_0\right),
\qquad x\longrightarrow\pm\infty,
\end{equation}
其中 $\Omega=-\alpha Q^2+\dfrac{2\eta_0}{a}$. 非零光场对应双叶谱曲面 $\rho^2=\lambda^2+Q^2$, 介质则在传播矩阵中引入实共振极点 $\lambda=a$.

\section{相关工作与主要结果}
Li、He 与 Porsezian\cite{LHP2013}给出了 HMB 的 Lax 对和 Darboux 变换, Zhao 与 Zhaqilao\cite{ZZ2026}研究了非零背景上的呼吸子和怪波. Li、Xia 与 Wei\cite{LXW2021}建立了零光场背景下的无反射 RHP. Liu 与 Liu\cite{LL2026}的非零背景施加于粒子数反转, 光场仍在无穷远趋于零; 条件 \eqref{eq:targetboundary} 则直接改变空间谱矩阵.

非零边界逆散射的谱分析与无反射重构可参见范恩贵\cite{Fan2022}及 Zhao--Fan\cite{ZhaoFan2019}. Zhang、Chen 与 Yan\cite{ZCY2020}讨论了 Hirota 方程的非零边界散射及相位约束, Biondini 等\cite{BGKL2019}建立了含频率相关介质的两能级系统的非零光场散射理论.

本文由背景谱问题建立 Jost 解析列、散射约化与演化, 构造满足两处归一化和四点留数条件的矩阵 RHP. 无反射时将其化为有限维系统, 利用 Schur 行列式的正因子分解直接求得纯虚一轨道的唯一显式解. 由无穷远和双共振重构三个场, 再通过矩阵对数导数恢复完整 Lax 对, 证明 HMB 方程及同相边界. 此外得到共振散射相容性与有限阶正则性判据, 并确定一轨道解的幅值界、局域化、传播速度和内部周期; 相应数值求值与动力学图像集中于最后一章.
